\documentclass[12pt]{article}
\usepackage[T1]{fontenc}
\usepackage[utf8]{inputenc}
\usepackage[polish]{babel}
\usepackage{lmodern}
\usepackage[a4paper,margin=2.4cm]{geometry}
\usepackage{xcolor}
\usepackage{graphicx}
\usepackage{amsmath,amssymb}
\usepackage{booktabs}
\usepackage{array}
\usepackage{enumitem}
\usepackage{hyperref}
\usepackage{tcolorbox}
\usepackage{fancyhdr}
\definecolor{WEIIPrimary}{RGB}{0,130,150}
\definecolor{WEIIPrimaryDark}{RGB}{0,86,100}
\definecolor{WEIIPrimaryLight}{RGB}{229,246,248}
\definecolor{PRzNavy}{RGB}{31,63,142}
\definecolor{RLTextDark}{RGB}{34,34,34}
\definecolor{RLGray}{RGB}{90,90,90}
\definecolor{RLLightGray}{RGB}{245,246,248}
\definecolor{RLWarning}{RGB}{184,86,0}
\definecolor{RLExample}{RGB}{52,112,64}
\definecolor{RLCodeBg}{RGB}{248,248,248}

\usepackage{listings}
\usepackage{xcolor}

\lstdefinestyle{rlpython}{%
    language=Python,
    basicstyle=\ttfamily\scriptsize,
    keywordstyle=\color{WEIIPrimaryDark}\bfseries,
    commentstyle=\color{RLGray}\itshape,
    stringstyle=\color{RLExample},
    numbers=left,
    numberstyle=\tiny\color{RLGray},
    numbersep=6pt,
    showstringspaces=false,
    breaklines=true,
    tabsize=4,
    frame=single,
    rulecolor=\color{WEIIPrimary!35},
    backgroundcolor=\color{RLCodeBg},
    captionpos=b
}

\lstset{style=rlpython}


\hypersetup{colorlinks=true,linkcolor=WEIIPrimaryDark,urlcolor=WEIIPrimaryDark,citecolor=WEIIPrimaryDark}
\setlength{\headheight}{15pt}
\pagestyle{fancy}
\fancyhf{}
\lhead{Uczenie się ze wzmocnieniem}
\rhead{WEiI PRz}
\cfoot{\thepage}

\newtcolorbox{taskbox}[1]{colback=WEIIPrimaryLight,colframe=WEIIPrimary,title=\textbf{#1}}
\newtcolorbox{notebox}[1]{colback=RLLightGray,colframe=RLGray,title=\textbf{#1}}
\newtcolorbox{warninghandout}[1]{colback=RLWarning!8,colframe=RLWarning,title=\textbf{#1}}

\newcommand{\handouttitle}[3]{%
\begin{center}
{\Large\bfseries #1}\\[0.2cm]
{\large #2}\\[0.15cm]
{\small #3}
\end{center}
\vspace{0.4cm}
\hrule
\vspace{0.6cm}
}

\newcommand{\EnableLabNumbering}{%
  \renewcommand{\thesection}{L\arabic{section}}%
  \renewcommand{\thesubsection}{\thesection.\arabic{subsection}}%
  \renewcommand{\thesubsubsection}{\thesubsection.\arabic{subsubsection}}%
}


\title{Raport projektowy\\Uczenie się ze wzmocnieniem}
\author{Piotr Janusz}
\date{2026/2027}

\begin{document}
\maketitle
\tableofcontents
\newpage

\section{Informacje podstawowe}

\begin{center}
\begin{tabular}{p{0.34\textwidth}p{0.56\textwidth}}
\toprule
Element & Informacja \\
\midrule
Imię i nazwisko & Piotr Janusz \\
Numer albumu & 179506 \\
Kierunek/specjalność & 3EF-DI/AI \\
Tytuł projektu & Czteronożny robot "Crawler" \\
Link do repozytorium & \url{https://github.com/pjanusz538/Crawler} \\
Data oddania & Do uzupełnienia \\
\bottomrule
\end{tabular}
\end{center}

\section{Temat i cel projektu}

\subsection{Opis problemu}
%Należy krótko opisać problem rozwiązywany w~projekcie. 
%Opis powinien być zrozumiały bez czytania kodu źródłowego.

Projekt dotyczy sterowania wirtualnym robotem czworonożnym w środowisku 
symulacyjnym Unity. Problem polega na odpowiednim
koordynowaniu ruchów kończyn, aby robot mógł zachować równowagę
i poruszać się w kontrolowany sposób.

\subsection{Cel agenta}
%Należy wyjaśnić, co agent ma osiągnąć oraz kiedy jego działanie można uznać za poprawne lub skuteczne.

Celem agenta jest nauczenie robota samodzielnego wstawania,
utrzymywania stabilnej pozycji oraz przemieszczania się
do wyznaczonego celu. Sukces oznacza zdolność do wykonywania
tych zadań bez utraty równowagi i przewrócenia się.

\subsection{Uzasadnienie zastosowania uczenia się ze wzmocnieniem}
%Należy wyjaśnić, dlaczego problem można rozpatrywać jako problem uczenia się ze wzmocnieniem. W szczególności należy wskazać:
%\begin{itemize}
%    \item jakie decyzje są podejmowane sekwencyjnie,
%    \item w jaki sposób decyzje wpływają na przyszły stan środowiska,
%    \item dlaczego istotny jest wynik długoterminowy, a nie tylko natychmiastowa nagroda,
%    \item dlaczego zastosowanie agenta uczącego się jest uzasadnione.
%\end{itemize}

Agent podejmuje sekwencyjne decyzje dotyczące ruchu stawów robota. 
Każdy ruch wpływa na jego położenie, równowagę i możliwości wykonania 
kolejnych działań. Istotne jest osiągnięcie długoterminowego celu, 
ponieważ pojedynczy ruch nie zawsze prowadzi do poprawnego poruszania się. 
Uczenie ze wzmocnieniem pozwala agentowi samodzielnie nauczyć się skutecznego 
sterowania poprzez interakcję ze środowiskiem i otrzymywane nagrody lub kary.

\section{Formalizacja problemu RL/MDP}

\subsection{Agent i środowisko}
%Należy opisać agenta, środowisko oraz sposób ich interakcji. Warto wskazać, co dzieje się w jednym kroku epizodu.

Agentem jest wirtualny robot czworonożny sterowany za pomocą
algorytmu PPO. Środowiskiem jest symulacja fizyczna
przygotowana w silniku Unity, zawierająca robota, podłoże
oraz wyznaczony cel.

W każdym kroku agent otrzymuje obserwacje opisujące aktualny
stan robota i położenie celu. Na ich podstawie wybiera akcje
sterujące przegubami kończyn. Następnie środowisko wykonuje
symulację fizyczną, aktualizuje położenie robota i oblicza
nagrodę za podjęte działanie.

\subsection{Stan lub obserwacja}

Agent nie obserwuje pełnego stanu środowiska, lecz otrzymuje
zestaw obserwacji opisujących robota i jego otoczenie.
Obejmują one prędkość liniową i kątową, orientację ciała,
położenie celu względem robota, aktualne obroty przegubów
oraz informację o kontakcie stóp z podłożem.

\begin{center}
\begin{tabular}{p{0.30\textwidth}p{0.60\textwidth}}
\toprule
Składowa stanu/obserwacji & Znaczenie \\
\midrule
Prędkość liniowa i kątowa & Informuje o ruchu i obrocie robota. \\
\midrule
Orientacja ciała & Określa ustawienie robota i jego odchylenie od pionu. \\
\midrule
Położenie celu & Określa kierunek i odległość do celu względem robota. \\
\midrule
Obroty przegubów & Opisują aktualne ustawienie kończyn robota. \\
\midrule
Kontakt stóp z podłożem & Informuje, które stopy mają kontakt z podłożem. \\
\bottomrule
\end{tabular}
\end{center}

\subsection{Akcje}

Agent korzysta z ciągłej przestrzeni akcji. Każda akcja przyjmuje
wartość z przedziału od $-1$ do $1$, która jest przeliczana
na docelowy kąt obrotu odpowiedniego przegubu. Sterowanie
odbywa się poprzez ustawianie docelowych rotacji przegubów
kończyn robota.

\begin{center}
\begin{tabular}{p{0.30\textwidth}p{0.60\textwidth}}
\toprule
Akcja & Opis \\
\midrule
Obrót górnych przegubów & Sterowanie obrotem w osiach X i Z. \\
\midrule
Obrót dolnych przegubów & Sterowanie obrotem w osi X. \\
\midrule
Wartość akcji & Wartość z przedziału $[-1,1]$, przeliczana na zakres kątów przegubu. \\
\bottomrule
\end{tabular}
\end{center}

Wszystkie opisane akcje są dostępne w każdym kroku sterowania. Ograniczenia fizyczne środowiska 
oraz symulacja w Unity zapewniają, że robot nie może wykonać niemożliwych ruchów np. obrót kończyn 
do wnętrza ciała czy przenikanie przez części ciała.

\subsection{Funkcja nagrody}

Funkcja nagrody służy do oceny działań agenta. Nagrody dodatnie
przyznawane są za zbliżanie się do celu, poprawę orientacji
i wysokości ciała oraz kontakt stóp z podłożem. Dodatkowa
nagroda jest przyznawana za utrzymanie stabilnej pozycji.

Przyjęta funkcja nagrody ma postać:

\begin{equation}
\begin{aligned}
r_t ={}&
0.5 \max(0, d_{t-1}-d_t) \\
&+ 0.001 \max(0, \mathbf{f}\cdot\mathbf{g}) \\
&+ 0.05(u_t-u_{t-1}) \\
&+ 0.1(h_t-h_{t-1}) \\
&+ 0.0005 n_t \\
&+ r_{\mathrm{stabilność}}
\end{aligned}
\end{equation}

gdzie:
\begin{itemize}
    \item $d_t$ oznacza odległość robota od celu,
    \item $\mathbf{f}$ oznacza kierunek przodu robota,
    \item $\mathbf{g}$ oznacza kierunek do celu,
    \item $u_t$ oznacza miarę pionowej orientacji robota,
    \item $h_t$ oznacza wysokość ciała robota,
    \item $n_t$ oznacza liczbę stóp mających kontakt z podłożem.
\end{itemize}

Nagroda za stabilność wynosi $0.002$, gdy robot jest odpowiednio
wyprostowany, osiąga wymaganą wysokość i utrzymuje kontakt
z podłożem wymaganą liczbą stóp. W przeciwnym przypadku
agent otrzymuje karę $-0.0005$.

Za osiągnięcie celu przyznawana jest nagroda $1.0$, natomiast
za przewrócenie się robota przyznawana jest kara $-1.0$.

Wartości współczynników określają znaczenie poszczególnych
składowych nagrody i pozwalają różnicować ich wpływ na proces
uczenia.

Dokładne wartości i dodatkowe kary/nagrody mogą ulec zmianie w trakcie 
eksperymentów w celu optymalizacji procesu uczenia i rozbudowy funkcji nagrody.

\subsection{Warunki zakończenia epizodu}

Epizod zostaje zakończony w przypadku przewrócenia robota,
czyli przekroczenia ustalonego progu odchylenia od pionu.
W takiej sytuacji agent otrzymuje karę $-1.0$.
Epizod kończy się również po osiągnięciu limitu
10\,000 kroków symulacji.

Osiągnięcie celu nie powoduje zakończenia epizodu.
Agent otrzymuje za to nagrodę dodatnią, a cel zostaje
wylosowany ponownie w innym miejscu.

\subsection{Założenia i ograniczenia środowiska}

Środowisko projektu zostało przygotowane w silniku Unity
z wykorzystaniem symulacji fizycznej. Robot porusza się
po płaskim podłożu, a jego ruch jest ograniczony do
możliwości fizycznych modelu oraz zakresów obrotu przegubów.

Środowisko nie uwzględnia obecnie nierówności terenu,
przeszkód ani zmiennych warunków otoczenia. Agent nie posiada
również dodatkowych czujników pozwalających wykrywać
przeszkody lub analizować ukształtowanie podłoża.

Ograniczeniem jest także uproszczony model fizyczny robota,
który może nie odwzorowywać w pełni zachowania rzeczywistej
konstrukcji. W przyszłości planowane jest rozszerzenie
środowiska o nierówny teren oraz mechanizmy wykrywania
i omijania przeszkód.

\section{Implementacja}

\subsection{Struktura repozytorium}

Należy krótko opisać strukturę repozytorium i najważniejsze pliki projektu.

\begin{verbatim}
repozytorium/
|-- src/
|-- experiments/
|-- report/
|-- results/
`-- README.md
\end{verbatim}

\subsection{Środowisko}

Należy opisać sposób implementacji środowiska, w szczególności funkcje inicjalizacji epizodu oraz wykonania kroku środowiska.

\subsection{Agent bazowy}

Należy opisać agenta losowego, heurystycznego albo inny baseline wykorzystany do porównania z agentem uczącym się.

\subsection{Agent uczący się}

Należy opisać zastosowany algorytm RL lub DRL, jego parametry oraz sposób implementacji. W przypadku Q-tab należy opisać reprezentację tablicy wartości. W przypadku DRL należy opisać wejścia modelu, wyjścia modelu, architekturę oraz sposób uczenia.

\section{Plan eksperymentów}

Należy opisać eksperymenty zaplanowane w celu oceny działania agenta.

\begin{center}
\begin{tabular}{p{0.24\textwidth}p{0.30\textwidth}p{0.30\textwidth}}
\toprule
Eksperyment & Co jest zmieniane? & Co jest mierzone? \\
\midrule
E1 & Do uzupełnienia & Do uzupełnienia \\
E2 & Do uzupełnienia & Do uzupełnienia \\
E3 & Do uzupełnienia & Do uzupełnienia \\
\bottomrule
\end{tabular}
\end{center}

W tej sekcji należy wskazać:
\begin{itemize}
    \item liczbę epizodów uczenia,
    \item liczbę powtórzeń eksperymentu albo uzasadnienie braku powtórzeń,
    \item parametry uczenia,
    \item metryki oceny,
    \item warianty środowiska, map lub funkcji nagrody,
    \item sposób porównania z agentem bazowym.
\end{itemize}

\section{Wyniki}

W tej sekcji należy przedstawić wyniki eksperymentów. Wyniki powinny być pokazane w postaci tabel, wykresów lub czytelnych zestawień.

\begin{center}
\begin{tabular}{p{0.28\textwidth}p{0.20\textwidth}p{0.20\textwidth}p{0.20\textwidth}}
\toprule
Wariant & Średni zwrot & Skuteczność & Długość epizodu \\
\midrule
Agent bazowy & Do uzupełnienia & Do uzupełnienia & Do uzupełnienia \\
Agent RL/DRL & Do uzupełnienia & Do uzupełnienia & Do uzupełnienia \\
\bottomrule
\end{tabular}
\end{center}

\begin{figure}[h]
    \centering
    \fbox{\parbox{0.82\textwidth}{\centering Miejsce na wykres procesu uczenia}}
    \caption{Przykładowy wykres procesu uczenia agenta.}
\end{figure}

\section{Analiza i wnioski}

Należy zinterpretować wyniki. W tej części należy odpowiedzieć m.in. na pytania:
\begin{itemize}
    \item Czy agent uczący się działa lepiej niż agent bazowy?
    \item Jak zmieniało się zachowanie agenta w trakcie uczenia?
    \item Które elementy środowiska były najtrudniejsze?
    \item Czy funkcja nagrody prowadziła do oczekiwanego zachowania?
    \item Czy wystąpiły przypadki niepożądanego zachowania agenta?
    \item Jakie są najważniejsze ograniczenia rozwiązania?
\end{itemize}

\section{Instrukcja uruchomienia}

Należy opisać sposób uruchomienia projektu od czystego repozytorium. Przykładowo:

\begin{verbatim}
git clone https://github.com/...
cd nazwa-repozytorium
python -m venv .venv
.venv\Scripts\activate
pip install -r requirements.txt
python src/main.py
\end{verbatim}

Należy również wskazać, jak odtworzyć eksperymenty oraz gdzie zapisywane są wyniki.

\section{Wsparcie AI}

W tej sekcji należy obowiązkowo opisać wykorzystanie narzędzi sztucznej inteligencji. Jeżeli student nie korzystał z takich narzędzi, należy wpisać jednoznaczną informację, że nie korzystano ze wsparcia AI.

\begin{center}
\begin{tabular}{p{0.13\textwidth}p{0.19\textwidth}p{0.21\textwidth}p{0.27\textwidth}}
\toprule
Data & Narzędzie/model & Cel użycia & Transkrypt / link / zakres \\
\midrule
Do uzupełnienia & Do uzupełnienia & Do uzupełnienia & Do uzupełnienia \\
Do uzupełnienia & Do uzupełnienia & Do uzupełnienia & Do uzupełnienia \\
\bottomrule
\end{tabular}
\end{center}

Należy zapisać transkrypty rozmów z modelami AI wraz z datą oraz nazwą modelu. Jeżeli rozmowa odbywała się poza narzędziami programistycznymi, należy dodatkowo udostępnić całą rozmowę poprzez link, o ile użyte narzędzie umożliwia takie udostępnienie.

Student powinien także wyjaśnić, które elementy wygenerowane lub zasugerowane przez AI zostały samodzielnie sprawdzone, zmodyfikowane albo odrzucone.

\section{Literatura}

W tej sekcji należy wskazać wykorzystane źródła: książki, artykuły naukowe, dokumentację bibliotek, materiały wykładowe oraz inne źródła użyte podczas pracy nad projektem.

\begin{thebibliography}{9}
\bibitem{sutton_barto}
R. S. Sutton, A. G. Barto, \textit{Reinforcement Learning: An Introduction}, 2nd ed., MIT Press, 2018.
\end{thebibliography}

\end{document}
